\scanpage{107}
\mspage{(10)}
\[
 R^o\subset S\subset\widehat{R^o}\tag*{(8.2.15)}
\]
\deleted{\unsure{donc}} et le foncteur
\[
 \unsure{\pi|\pi_{\xi}}:R\times S\longrightarrow(\mathrm{Ens})
\]
est donné par
\[
 \pi(c,\xi)=\deleted{\unsure{\mathrm{Hom}_{S}(c}}\ \xi(c)=\mathrm{Hom}_{\widehat{R^o}}(c,\xi)=\mathrm{Hom}_{S}(c,\xi)\tag*{(8.2.16)}
\]
i.e. \deleted{\unsure{le composé}}\inserted{l'opposé} du foncteur (8.2.13) est \deleted{induit par le}
\deleted{foncteur} composé des deux foncteurs pleinement fidèles
\[
 R^{\deleted{\unsure{o}}}\underset{(8.2.15)}{\lhook\joinrel\longrightarrow}S^o\lhook\joinrel\longrightarrow\widehat{S^o}=\underline{\mathrm{Hom}}\,(S,(\mathrm{Ens}))\tag*{(8.2.17)}
\]
donc est pl. fidèle. Ainsi $R$ est connu \inserted{à équiv. près} quand
on connaît son image dans $\widehat{S^o}$. \deleted{\unsure{En} \unsure{termes}}.

Pour résumer, si $T$ est un $\lambda$-type $(\underset{\rightarrow}{\lambda}=\varnothing)$, alors
\deleted{pour \unsure{toute} $C\in\mathrm{Ob}\,\mathrm{Cat}_{\lambda}$,} $S=T(\mathrm{Ens})$ est muni d'une sous-catégorie
pleine $\Sigma\subset S$, formée des $\varphi\in\mathrm{Ob}\,S$ tels que pour
tout $C\in\mathrm{Ob}\,\mathrm{Cat}_{\lambda}\deleted{\unsure{C}}$ et tout \deleted{\unsure{fon}} $\xi\in C_\lambda$, \deleted{\unsure{on a un fon}}
\deleted{\unsure{s'expr}} définissant un \deleted{\unsure{hom}} foncteur
\[
 \widetilde\xi:C^o\longrightarrow S\qquad X\longmapsto\mathrm{Hom}(X,\xi)
\]
le foncteur composé avec $\mathrm{Hom}_S(\varphi,-)$ soit
représentable\unsure{s} :
\[
 X\longrightarrow\mathrm{Hom}_S(\varphi,\deleted{\unsure{\widetilde\xi(X)}}\mathrm{Hom}(X,\xi))
\]
$\Sigma$ est stable dans $S$ par $\varinjlim$ \inserted{(\uline{inductives})} de type $\in\lambda$ \deleted{\unsure{\xi}} [car $C$
est stable dans $\widehat C$ par $\varprojlim$ de type $\lambda$, $\forall C\in\mathrm{Ob}\,\mathrm{Cat}_{\lambda}$]
et pour tte $C$ comme dessus, on a \unsure{deux|des} foncteurs
\[
 \Sigma^o\xrightarrow{\inserted{\lambda\text{-foncteur}}}\underline{\mathrm{Hom}}\,(T(C),C),\qquad T(C)\longrightarrow\underline{\mathrm{Hom}}_\lambda(\Sigma^o,C)
\]
On voit alors que l'on trouve ainsi une
% archive sheet number: 106
